% page 438 of the facsimile (image p-437 of the scan)
% block 157.5 x 246.3 mm ; left margin 8.0 mm
\pgc{-12.700mm}{0.000mm}{
\l{\cel{3}430}
\l{}
\l{\sou{pelagika}.\cel{2}(\sou{zool}.) Qua vivas en maro larja, litoro-fore.DEFIS.}
\l{}
\l{\sou{pelagro}. (\sou{patol.})\cel{3}Morbo di pelo, qua divenas kakexio,}
\l{\cel{3}e di qua la kauzo esas la indijo di vitamini en la}
\l{\cel{3}nutrivi. - DEFIS.}
\l{}
\l{\sou{pelargonio}. (\sou{bot.}) Genero de \textquotedbl{}geraniacei\textquotedbl{}. - L.\cel{1}\sou{pelargo-}}
\l{\cel{3}\sou{nium.}\cel{3}DEFIS.}
\l{}
\l{\sou{peleri}no. Granda kolumo jacanta (manteleto) qua tegas la}
\l{\cel{3}shultro e la pektoro. - DEFIR.}
\l{}
\l{\sou{pelikano}. (\sou{zool.)}\cel{2}Ucelo aquala, kun beko larja, di qua la}
\l{\cel{3}mandibulo esas garnisita ye granda posho membrana aden}
\l{\cel{3}qua lu rezervas sua peskajo. - L.\cel{1}\sou{pelecanus}.- DEFIRS.}
\l{}
\l{\sou{peliku}lo. I. Pelo, membrano, tre dina. - II. Membrano tenua,}
\l{\cel{3}dina, qua formacesas sur la surfaco di ula substanci (je-}
\l{\cel{3}lei, lakto boliita). - EFIS.}
\l{}
\l{\sou{peliso}. Mantelo vatizita. - EFIS.}
\l{}
\l{\sou{pelmelo}. Desordino maxima, qua memorigas kaoso. - EF.}
\l{}
\l{\sou{pelorio}. (\sou{bot.}) Anomalajo qua konsistas ye to, ke korolo}
\l{\cel{3}normala ne-segunregula divenas segun-regula. (Digitalo}
\l{\cel{3}esas exemplo pri lo).- DEFIS.}
\l{}
\l{\sou{peloto}.Mikra amaso qua havas la formo di baloneto. (Peloto}
\l{\cel{3}de nivofloki, de lino-filo, de lano-filo). - FS.}
\l{}
\l{\sou{pelotono}. (\sou{armeo)}.\cel{2}Mikra trupo de soldati, e, partikulare,}
\l{\cel{3}subdividuro di kompanio, di eskadrono. - DEFIS.}
\l{}
\l{\sou{pelvo.}\cel{2}(\sou{anat.}) Quaza zono osta, qua konstitucas la bazo di}
\l{\cel{3}la torso, e konsistas ye la du osti hanchala. - EFIS.}
\l{}
\l{\sou{pemikano.}\cel{2}Preparajo de karno sikigita e kondensita. - DEFIS}
\l{}
\l{\sou{penar}. (\sou{netrans}.)Facar granda esforco qua fatigas multe.-}
\l{\cel{3}EFIS.}
\l{}
\l{\sou{penacho}. Fasko de plumi, (maxim-multa-kaze) koloroza,\cel{2}inter-}
\l{\cel{3}klemita infre, fluktuanta supre. - EFIS.}
\l{}
\l{\sou{penato.}\cel{2}(en\cel{1}\sou{Roma, en la epoki antiqua}). Singla de la deaji}
\l{\cel{3}qui protektis la hemo, e quin onu kultas che su. - DEFIRS}
\l{}
\l{\sou{pendar}. I. (\sou{trans}.) Ligar per la parto supra, diste de sulo.}
\l{\cel{3}II. (\sou{netrans}.) Ligar a jibeto, per kordo, cirkondante la}
\l{\cel{3}kolo por mortigar per strangulo. III. (\sou{netrans}.)Esar}
\l{\cel{3}ligita per la parto supra, diste de sulo. - EFS.}
\l{}
\l{\sou{pendentivo}.\cel{2}(\sou{arkitekt.}) Parto de vulto, inter la granda arki}
\l{\cel{3}qui suportas kupolo. - DEF.}
}
